%%%%%%%%%%%%%%%%%%%%%%%%%%%%%%%%%%%%%%%%%
% GWO Designer - Overview document
% A short, simplified companion to the detailed report (report.tex):
% intro + paper reference, recovery QR/link, the four main result plots
% with captions, and the full list of input parameters.
% Uses the SelfArx class (see report.tex) for a consistent look.
%%%%%%%%%%%%%%%%%%%%%%%%%%%%%%%%%%%%%%%%%

\documentclass[fleqn,10pt]{<approot>/designer/latexelements/SelfArx}

\usepackage[english]{babel}
\selectlanguage{english}

\DeclareCaptionFont{blue}{\color{color3}}
\usepackage[font={blue,footnotesize,sf}, labelfont={blue,footnotesize,sf,bf}, justification=justified]{caption}

\usepackage{upgreek}
\usepackage{bm}
\usepackage{textcomp}
\usepackage[autolanguage]{numprint}
\usepackage{expl3,siunitx}
\usepackage{booktabs}
\usepackage{longtable}
\usepackage[table]{xcolor}
\usepackage{verbatimbox}
% verbatimbox/readarray version skew in modern TeX Live: \addvbuffer -> undefined
% \getargsC. \addvbuffer only adds cosmetic padding, so make it a passthrough.
\makeatletter\renewcommand{\addvbuffer}[2][]{#2}\makeatother

\setlength{\columnsep}{0.55cm}
\setlength{\fboxrule}{0.75pt}

\definecolor{color3}{RGB}{60,60,60}
\definecolor{color1}{RGB}{0,60,80}
\definecolor{color2}{RGB}{70,130,140}

\usepackage{hyperref}
\hypersetup{hidelinks,colorlinks,breaklinks=true,urlcolor=color1,citecolor=color1,linkcolor=color1,bookmarksopen=false,pdftitle={GWO Designer - Overview},pdfauthor={spacegravity.org/designer}}

\JournalInfo{Auto-generated overview \textbullet\ \today}
\Archive{spacegravity.org/designer \textbullet\ recovery code: <session>}
\PaperTitle{The <detectorName> Observatory}
\PaperSubTitle{Design overview of a spaceborne gravitational wave detector}
\Authors{<author>}
\affiliation{\textit{<affiliation>}}
\affiliation{Generated with the Gravitational Wave Observatory Designer, \href{http://spacegravity.org/designer}{spacegravity.org/designer}}
\Keywords{Gravitational Waves --- Space Interferometer --- LISA}
\newcommand{\keywordname}{Keywords}

\Abstract{This is a short overview of the \textbf{<detectorName>} gravitational wave observatory, a <numberOfLinks>-link configuration with an interferometric arm length of <armLength>\,<armLengthUnit>, designed with the online Gravitational Wave Observatory Designer. It collects the chosen input parameters, the resulting sensitivity plots, and a recovery code for reopening the design. A full technical report with all underlying calculations is available separately from the tool.}

\begin{document}

\flushbottom
\maketitle
\thispagestyle{empty}

%----------------------------------------------------------------------------------------
\section*{About this document}
The plots and parameters below describe one observatory design produced with the
Gravitational Wave Observatory Designer. The method, noise models, and full derivations
are described in the related publication: S.~Barke \textit{et al.}, ``Towards a
gravitational wave observatory designer: sensitivity limits of spaceborne detectors'',
\textit{Class.\ Quantum Grav.} \textbf{32}, 095004 (2015). Results are provided for
guidance and comparison of relative parametric dependencies; verify them independently
before relying on them.

\paragraph{Recovery.} This design is stored online under the twelve-digit recovery code
\textbf{\texttt{<session>}}. To reopen it with all parameters, visit
\href{http://spacegravity.org/designer/\#rc=<session>}{spacegravity.org/designer} and
enter the code, use the permalink
\href{http://spacegravity.org/designer/\#rc=<session>}{http://spacegravity.org/designer/\#rc=<session>},
or scan the code below.

\begin{center}\includegraphics[width=0.28\linewidth]{<qrplot>}\end{center}

%----------------------------------------------------------------------------------------
\begin{figure*}[!htb]\centering
\includegraphics[width=0.92\linewidth]{<overviewplot>}
\caption{\textbf{Mission overview.} Comparison of the strain sensitivity of the designed
observatory with reference detector concepts. All submitted parameter sets are shown.}
\end{figure*}

\begin{figure*}[!htb]\centering
\includegraphics[width=0.92\linewidth]{<displacementplot>}
\caption{\textbf{Displacement noise contributions.} Gravitational waves alter the distance
between spacecraft, producing a phase shift in the interferometric read-out. Several noise
sources are indistinguishable from a genuine spacecraft displacement and limit the
observatory's sensitivity; the considered contributions are shown over gravitational wave
frequency for the chosen parameters.}
\end{figure*}

\begin{figure*}[!htb]\centering
\includegraphics[width=0.92\linewidth]{<singleplot>}
\caption{\textbf{Single-link strain sensitivity.} Impact of gravitational waves on one laser
link, limited by the displacement noise budget and averaged over all sky positions and both
polarizations. The carrier read-out noise and the proof-mass acceleration noise are shown
separately.}
\end{figure*}

\begin{figure*}[!htb]\centering
\includegraphics[width=0.92\linewidth]{<fullplot>}
\caption{\textbf{Observatory and astrophysical sources.} Detection limit (signal-to-noise
ratio of one) of the full observatory, compared with characteristic strain amplitudes of
massive black-hole binaries, extreme mass-ratio inspirals (EMRIs), and ultra-compact
binaries.}
\end{figure*}

%----------------------------------------------------------------------------------------
\clearpage
\onecolumn
\section*{Input parameters}
\renewcommand{\arraystretch}{1.15}
\footnotesize
\begin{longtable}{@{}p{0.56\textwidth}p{0.22\textwidth}p{0.18\textwidth}@{}}
\toprule
\textbf{Parameter} & \textbf{Value} & \textbf{Unit} \\
\midrule
\endfirsthead
\toprule
\textbf{Parameter} & \textbf{Value} & \textbf{Unit} \\
\midrule
\endhead
\bottomrule
\endfoot
<parametertable>
\end{longtable}

\end{document}
